\documentclass[a4paper,11pt]{book}
\usepackage[utf8x]{inputenc}

\usepackage{graphicx}

%% or use the epsfig package if you prefer to use the old commands
%% \usepackage{epsfig}

%% The amssymb package provides various useful mathematical symbols
%% \usepackage{amssymb}
%% The amsthm package provides extended theorem environments
%% \usepackage{amsthm}

%% The lineno packages adds line numbers. Start line numbering with
%% \begin{linenumbers}, end it with \end{linenumbers}. Or switch it on
%% for the whole article with \linenumbers after \end{frontmatter}.
%% \usepackage{lineno}

%% \usepackage{placeins}

\usepackage{listings}
\usepackage{subfigure}
\usepackage{courier}
\usepackage{tikz}
\usetikzlibrary{shapes,positioning}
\usepackage{mdwlist}
\usepackage{hyperref}
\hypersetup{
    colorlinks,
    citecolor=black,
    filecolor=black,
    linkcolor=black,
    urlcolor=black
}
\usepackage{url}

\usepackage{color}

\definecolor{lightlightgray}{gray}{0.97}
\definecolor{lightlightgreen}{rgb}{0.8,1,0.8}
\definecolor{darkblue}{rgb}{0,0,0.8}
\definecolor{darkgreen}{rgb}{0,0.5,0}
\definecolor{LightSkyBlue}{rgb}{0.8,1,0.8}
\definecolor{MidnightBlue}{rgb}{0.8,1,0.8}
\definecolor{SandiaGray}{RGB}{130, 120, 111}
\definecolor{SandiaBlue}{RGB}{0, 50, 90}
\definecolor{SandiaLightBlue}{RGB}{0,150,180}
\definecolor{SandiaLightLightBlue}{RGB}{0,200,220}
\definecolor{SandiaLightLightLightBlue}{RGB}{0,228,238}
\definecolor{SandiaRed}{RGB}{130, 36, 51}
\lstset{language=C++,
           frame=single,
           basicstyle=\ttfamily\scriptsize,
           %%basicstyle=\ttfamily,
           %%basicstyle=\sffamily\normalsize,
           keywordstyle=\color{darkblue},
           backgroundcolor=\color{lightlightgray},
           commentstyle=\color{darkgreen}
           }

\tikzset{HighLightBox/.style = {draw=SandiaBlue, fill=SandiaLightLightLightBlue, line width=2pt,
                          rectangle, rounded corners, inner sep=10pt, inner ysep=20pt},
         fancytitle/.style ={fill=SandiaBlue, text=white, rounded corners}
         }



\begin{document}
\begin{minipage}[h]{\paperwidth}
  {\begin{tikzpicture}[remember picture,overlay]
    \node[] at (current page.north west)
      {\begin{tikzpicture}[remember picture, overlay, scale=0.5]
          \draw[fill=SandiaLightLightBlue,color=SandiaLightLightBlue] (0,-10.5cm) rectangle (8cm,5.5cm);
          \draw[fill=SandiaLightBlue,color=SandiaLightBlue] (0,-8.5cm) rectangle (6cm,3.5cm);
          \draw[fill=SandiaBlue,color=SandiaBlue] (0,-6.5cm) rectangle (4cm,1.5cm);          
          \draw[fill=SandiaLightBlue,color=SandiaLightBlue] (4cm,-6.5cm) circle [radius=4cm];
          \draw[fill=SandiaLightLightBlue,color=SandiaLightLightBlue] (6cm,-8.5cm) circle [radius=6cm];
          \draw[fill=white,color=white] (8cm,-10.5cm) circle [radius=8cm];
          \draw[fill=SandiaGray,color=SandiaGray] (0,0cm) rectangle (2\paperwidth,-0.5cm);
          \draw[fill=SandiaBlue,color=SandiaBlue] (0,-0.5cm) rectangle (2\paperwidth,-2.5cm);
       \end{tikzpicture}
      };
     \node[anchor=north west,xshift=-.05\paperwidth,rectangle,
               yshift=3.75cm, rounded corners=10pt,inner sep=5pt,
              fill=SandiaBlue]
              {\color{white} \LARGE Kokkos: Quick Reference};
   \end{tikzpicture}
   }

  {\begin{tikzpicture}[remember picture,overlay]
      \node[] at (current page.south east)
      {\begin{tikzpicture}[remember picture, overlay,scale=0.5]
          \draw[fill=SandiaLightLightBlue,color=SandiaLightLightBlue] (-8cm,11.5cm) rectangle (8cm,5.5cm);
          \draw[fill=SandiaLightBlue,color=SandiaLightBlue] (-6cm,9.5cm) rectangle (6cm,3.5cm);
          \draw[fill=SandiaBlue,color=SandiaBlue] (0,7.5cm) rectangle (-4cm,1.5cm);          
          \draw[fill=SandiaLightBlue,color=SandiaLightBlue] (-4cm,7.5cm) circle [radius=4cm];
          \draw[fill=SandiaLightLightBlue,color=SandiaLightLightBlue] (-6cm,9.3cm) circle [radius=6cm];
          \draw[fill=white,color=white] (-8cm,11.5cm) circle [radius=8cm];
        \draw[fill=SandiaBlue,color=SandiaBlue] (-2\paperwidth,0.5cm) rectangle
          (\paperwidth,3.5cm);
        \draw[fill=SandiaGray,color=SandiaGray] (-2\paperwidth,0) rectangle
          (\paperwidth,0.5cm);
       \end{tikzpicture}
      };                
      \node[xshift=0.75\paperwidth,rectangle, yshift=2.65cm] at (current page.south west)
              {\includegraphics[height=0.3in]{figures/SNL_Stacked_Black_Blue.jpg}};
      \node[xshift=0.6\paperwidth,rectangle, yshift=2.65cm] at (current page.south west)
              {\includegraphics[height=0.45in]{figures/DOEbwlogo.pdf}};
           
   \end{tikzpicture}
   }

  {\begin{tikzpicture}[remember picture,overlay,trim left]
     \node[HighLightBox,yshift=-12cm,xshift=0.265\paperwidth] at (current page.north west) (ParallelPatterns)
     {\begin{minipage}[h]{0.4\paperwidth}
     \begin{small}
Algorithms are expressed in parallel patterns for all which applies it:
\begin{itemize}\itemsep1pt \parskip0pt \parsep0pt
\item execute every element of index set
\item order and concurrency are undefined
\item \lstinline|MemberType| is used for \lstinline|ExecPolicy::member_type|
\item \lstinline|TagType| is used for \lstinline|ExecPolicy::tag_type|
\item \lstinline|TagType| is used if it is not \lstinline|void|
\end{itemize}
Using appropriate Exececution Policies parallel patterns can be nested.
\begin{lstlisting}
parallel_for( ExecPolicy , Functor);
\end{lstlisting}
\vspace{-0.4cm}
\begin{itemize}\itemsep1pt \parskip0pt \parsep0pt
\item calls \lstinline|Functor::operator([TagType,] MemberType)|
\end{itemize}
\begin{lstlisting}
parallel_reduce( ExecPolicy , Functor);
parallel_reduce( ExecPolicy , Functor, Result);
\end{lstlisting}
\vspace{-0.4cm}
\begin{itemize}\itemsep1pt \parskip0pt \parsep0pt
\item calls \lstinline|Functor::operator([TagType,] MemberType,| \lstinline|ValueType&)|
\item performs a deterministic reduction over \lstinline|ValueType|
\item \lstinline|ValueType| is inferred from \lstinline|Functor::operator|
\item \lstinline|Result| must be a 0D \lstinline|View| of \lstinline|ValueType| or
an lvalue of \lstinline|ValueType|
\item thread contributions are combined by \lstinline|Functor::join([TagType, ]volatile ValueType&,| 
 \lstinline|volatile ValueType&)| or \lstinline|ValueType::operator+=|  \lstinline|(volatile ValueType&)|
if \lstinline|join| is not defined
\item thread contributions are initialized by \lstinline|Functor::init([TagType, ] ValueType& )| 
or \lstinline|ValueType()| if \lstinline|init| is not defined
\item \lstinline|Functor::finalize([TagType, ] ValueType&)| is called for the final reduction
value if defined 
\end{itemize}
\begin{lstlisting}
parallel_scan( ExecPolicy , Functor);
\end{lstlisting}
\vspace{-0.4cm}
\begin{itemize}\itemsep1pt \parskip0pt \parsep0pt
\item calls \lstinline|Functor::operator([TagType,] MemberType,| \lstinline|ValueType&, bool final)|
\item if \lstinline|final == true| \lstinline|ValueType| is the prefix sum
\item the contributions have to be added to \lstinline|ValueType| irrespective of \lstinline|final|  
\end{itemize}
     \end{small}     
      \end{minipage}
      };
      \node[fancytitle, right=10pt] at (ParallelPatterns.north west) {Parallel Patterns};   
      \end{tikzpicture}
  }
  
  
  
  {\begin{tikzpicture}[remember picture,overlay]
     \node[HighLightBox,yshift=-14cm,xshift=0.75\paperwidth] at (current page.north west) (ExecutionPolicies)
     {\begin{minipage}[h]{0.4\paperwidth}
\begin{small}
Execution Policies govern how a function is executed and what the index space is over which
the parallel iteration is performed. 
\begin{lstlisting}
RangePolicy<[ExecSpace,] [TagType]>
      (begin, end)
RangePolicy<[ExecSpace,] [TagType,] INT>
      (begin, end)
RangePolicy<[ExecSpace,] [TagType,] IntType>
      (begin, end)
\end{lstlisting}
\vspace{-0.4cm}
\begin{itemize}\itemsep1pt \parskip0pt \parsep0pt
\item top level policy (i.e. not nested)
\item statically splits \lstinline|(begin, end]| interval over compute resources
\item \lstinline|ExecSpace| defaults to \lstinline|DefaultExecutionSpace|
\item \lstinline|TagType| defaults to \lstinline|void|
\item providing an integer number specifies chunk size for static splitting
\item providing an integer type sets the type of the loop index
\end{itemize}
\begin{lstlisting}
TeamPolicy<[ExecSpace,] [TagType]>
      (league_size, team_size[, vector_length])
TeamPolicy<[ExecSpace,] [TagType,] INT>
      (league_size, team_size[, vector_length])
TeamPolicy<[ExecSpace,] [TagType,] IntType>
      (league_size, team_size[, vector_length])
\end{lstlisting}
\vspace{-0.4cm}
\begin{itemize}\itemsep1pt \parskip0pt \parsep0pt
\item top level policy (i.e. not nested)
\item statically splits \lstinline|league_size| range over interval over compute resources
\item \lstinline|ExecSpace| defaults to \lstinline|DefaultExecutionSpace|
\item \lstinline|TagType| defaults to \lstinline|void|
\item providing an integer number specifies chunk size for static splitting
\item providing an integer type sets the type of the loop index
\item starts 2D or 3D parallelism: a league of thread teams with vector length
\end{itemize}
\begin{lstlisting}
TeamThreadRange(TeamPolicy<...>::member_type, 
                begin, end)
\end{lstlisting}
\vspace{-0.4cm}
\begin{itemize}\itemsep1pt \parskip0pt \parsep0pt
\item for each calling team split the range \lstinline|[begin,end)| over threads in the team
\item must be nested in call with \lstinline|TeamPolicy|
\item team collective
\end{itemize}        
\begin{lstlisting}
ThreadVectorRange(TeamPolicy<...>::member_type, 
                begin, end)
\end{lstlisting}
\vspace{-0.4cm}
\begin{itemize}\itemsep1pt \parskip0pt \parsep0pt
\item for each calling thread split the range \lstinline|[begin,end)| over vector lanes
\item must be nested in call with \lstinline|TeamPolicy|
\end{itemize}        
\end{small}
      \end{minipage}
      };
      \node[fancytitle, right=10pt] at (ExecutionPolicies.north west) {Execution Policies};   
      \end{tikzpicture}
  }

  {
  \begin{tikzpicture}[remember picture,overlay,trim left]
     \node[HighLightBox,yshift=-24.75cm,xshift=0.265\paperwidth] at (current page.north west) (ExecutionSpace)
     {\begin{minipage}[h]{0.4\paperwidth}
     \begin{small}
\begin{lstlisting}
Serial
\end{lstlisting}
\vspace{-0.4cm}
\begin{itemize}\itemsep1pt \parskip0pt \parsep0pt
\item basic serial execution
\end{itemize}
\vspace{-0.4cm}
\begin{lstlisting}
OpenMP
\end{lstlisting}
\vspace{-0.4cm}
\begin{itemize}\itemsep1pt \parskip0pt \parsep0pt
\item using OpenMP 3.1 as backend
\end{itemize}
\vspace{-0.4cm}
\begin{lstlisting}
Threads
\end{lstlisting}
\vspace{-0.4cm}
\begin{itemize}\itemsep1pt \parskip0pt \parsep0pt
\item using Pthreads on Unix based systems
\end{itemize}
\vspace{-0.4cm}
\begin{lstlisting}
Cuda
\end{lstlisting}
\vspace{-0.4cm}
\begin{itemize}\itemsep1pt \parskip0pt \parsep0pt
\item using NVIDIAs Cuda 6.5 or higher as backend
\end{itemize}
\vspace{-0.4cm}
\end{small}
      \end{minipage}
      };
      \node[fancytitle, right=10pt] at (ExecutionSpace.north west) {Execution Spaces};   
   \end{tikzpicture}
  }

\end{minipage}
\newpage
\begin{minipage}[h]{\paperwidth}
  {\begin{tikzpicture}[remember picture,overlay]
    \node[] at (current page.north west)
      {\begin{tikzpicture}[remember picture, overlay, scale=0.5]
          \draw[fill=SandiaLightLightBlue,color=SandiaLightLightBlue] (0,-10.5cm) rectangle (8cm,5.5cm);
          \draw[fill=SandiaLightBlue,color=SandiaLightBlue] (0,-8.5cm) rectangle (6cm,3.5cm);
          \draw[fill=SandiaBlue,color=SandiaBlue] (0,-6.5cm) rectangle (4cm,1.5cm);          
          \draw[fill=SandiaLightBlue,color=SandiaLightBlue] (4cm,-6.5cm) circle [radius=4cm];
          \draw[fill=SandiaLightLightBlue,color=SandiaLightLightBlue] (6cm,-8.5cm) circle [radius=6cm];
          \draw[fill=white,color=white] (8cm,-10.5cm) circle [radius=8cm];
          \draw[fill=SandiaGray,color=SandiaGray] (0,0cm) rectangle (2\paperwidth,-0.5cm);
          \draw[fill=SandiaBlue,color=SandiaBlue] (0,-0.5cm) rectangle (2\paperwidth,-2.5cm);
       \end{tikzpicture}
      };
     \node[anchor=north west,xshift=-.05\paperwidth,rectangle,
               yshift=3.75cm, rounded corners=10pt,inner sep=5pt,
              fill=SandiaBlue]
              {\color{white} \LARGE Kokkos: Quick Reference};
   \end{tikzpicture}
   }

  {\begin{tikzpicture}[remember picture,overlay]
      \node[] at (current page.south east)
      {\begin{tikzpicture}[remember picture, overlay,scale=0.5]
          \draw[fill=SandiaLightLightBlue,color=SandiaLightLightBlue] (-8cm,11.5cm) rectangle (8cm,5.5cm);
          \draw[fill=SandiaLightBlue,color=SandiaLightBlue] (-6cm,9.5cm) rectangle (6cm,3.5cm);
          \draw[fill=SandiaBlue,color=SandiaBlue] (0,7.5cm) rectangle (-4cm,1.5cm);          
          \draw[fill=SandiaLightBlue,color=SandiaLightBlue] (-4cm,7.5cm) circle [radius=4cm];
          \draw[fill=SandiaLightLightBlue,color=SandiaLightLightBlue] (-6cm,9.3cm) circle [radius=6cm];
          \draw[fill=white,color=white] (-8cm,11.5cm) circle [radius=8cm];
        \draw[fill=SandiaBlue,color=SandiaBlue] (-2\paperwidth,0.5cm) rectangle
          (\paperwidth,3.5cm);
        \draw[fill=SandiaGray,color=SandiaGray] (-2\paperwidth,0) rectangle
          (\paperwidth,0.5cm);
       \end{tikzpicture}
      };                
      \node[xshift=0.75\paperwidth,rectangle, yshift=2.65cm] at (current page.south west)
              {\includegraphics[height=0.3in]{figures/SNL_Stacked_Black_Blue.jpg}};
      \node[xshift=0.6\paperwidth,rectangle, yshift=2.65cm] at (current page.south west)
              {\includegraphics[height=0.45in]{figures/DOEbwlogo.pdf}};
           
   \end{tikzpicture}
   }

  {\begin{tikzpicture}[remember picture,overlay,trim left]
     \node[HighLightBox,yshift=-12cm,xshift=0.265\paperwidth] at (current page.north west) (ParallelPatterns)
     {\begin{minipage}[h]{0.4\paperwidth}
     \begin{small}
Algorithms are expressed in parallel patterns for all which applies it:
\begin{itemize}\itemsep1pt \parskip0pt \parsep0pt
\item execute every element of index set
\item order and concurrency are undefined
\item \lstinline|MemberType| is used for \lstinline|ExecPolicy::member_type|
\item \lstinline|TagType| is used for \lstinline|ExecPolicy::tag_type|
\item \lstinline|TagType| is used if it is not \lstinline|void|
\end{itemize}
Using appropriate Exececution Policies parallel patterns can be nested.
\begin{lstlisting}
parallel_for( ExecPolicy , Functor);
\end{lstlisting}
\vspace{-0.4cm}
\begin{itemize}\itemsep1pt \parskip0pt \parsep0pt
\item calls \lstinline|Functor::operator([TagType,] MemberType)|
\end{itemize}
\begin{lstlisting}
parallel_reduce( ExecPolicy , Functor);
parallel_reduce( ExecPolicy , Functor, Result);
\end{lstlisting}
\vspace{-0.4cm}
\begin{itemize}\itemsep1pt \parskip0pt \parsep0pt
\item calls \lstinline|Functor::operator([TagType,] MemberType,| \lstinline|ValueType&)|
\item performs a deterministic reduction over \lstinline|ValueType|
\item \lstinline|ValueType| is inferred from \lstinline|Functor::operator|
\item \lstinline|Result| must be a 0D \lstinline|View| of \lstinline|ValueType| or
an lvalue of \lstinline|ValueType|
\item thread contributions are combined by \lstinline|Functor::join([TagType, ]volatile ValueType&,| 
 \lstinline|volatile ValueType&)| or \lstinline|ValueType::operator+=|  \lstinline|(volatile ValueType&)|
if \lstinline|join| is not defined
\item thread contributions are initialized by \lstinline|Functor::init([TagType, ] ValueType& )| 
or \lstinline|ValueType()| if \lstinline|init| is not defined
\item \lstinline|Functor::finalize([TagType, ] ValueType&)| is called for the final reduction
value if defined 
\end{itemize}
\begin{lstlisting}
parallel_scan( ExecPolicy , Functor);
\end{lstlisting}
\vspace{-0.4cm}
\begin{itemize}\itemsep1pt \parskip0pt \parsep0pt
\item calls \lstinline|Functor::operator([TagType,] MemberType,| \lstinline|ValueType&, bool final)|
\item if \lstinline|final == true| \lstinline|ValueType| is the prefix sum
\item the contributions have to be added to \lstinline|ValueType| irrespective of \lstinline|final|  
\end{itemize}
     \end{small}     
      \end{minipage}
      };
      \node[fancytitle, right=10pt] at (ParallelPatterns.north west) {Parallel Patterns};   
      \end{tikzpicture}
  }
  
  
  
  {\begin{tikzpicture}[remember picture,overlay]
     \node[HighLightBox,yshift=-14cm,xshift=0.75\paperwidth] at (current page.north west) (ExecutionPolicies)
     {\begin{minipage}[h]{0.4\paperwidth}
\begin{small}
Execution Policies govern how a function is executed and what the index space is over which
the parallel iteration is performed. 
\begin{lstlisting}
RangePolicy<[ExecSpace,] [TagType]>
      (begin, end)
RangePolicy<[ExecSpace,] [TagType,] INT>
      (begin, end)
RangePolicy<[ExecSpace,] [TagType,] IntType>
      (begin, end)
\end{lstlisting}
\vspace{-0.4cm}
\begin{itemize}\itemsep1pt \parskip0pt \parsep0pt
\item top level policy (i.e. not nested)
\item statically splits \lstinline|(begin, end]| interval over compute resources
\item \lstinline|ExecSpace| defaults to \lstinline|DefaultExecutionSpace|
\item \lstinline|TagType| defaults to \lstinline|void|
\item providing an integer number specifies chunk size for static splitting
\item providing an integer type sets the type of the loop index
\end{itemize}
\begin{lstlisting}
TeamPolicy<[ExecSpace,] [TagType]>
      (league_size, team_size[, vector_length])
TeamPolicy<[ExecSpace,] [TagType,] INT>
      (league_size, team_size[, vector_length])
TeamPolicy<[ExecSpace,] [TagType,] IntType>
      (league_size, team_size[, vector_length])
\end{lstlisting}
\vspace{-0.4cm}
\begin{itemize}\itemsep1pt \parskip0pt \parsep0pt
\item top level policy (i.e. not nested)
\item statically splits \lstinline|league_size| range over interval over compute resources
\item \lstinline|ExecSpace| defaults to \lstinline|DefaultExecutionSpace|
\item \lstinline|TagType| defaults to \lstinline|void|
\item providing an integer number specifies chunk size for static splitting
\item providing an integer type sets the type of the loop index
\item starts 2D or 3D parallelism: a league of thread teams with vector length
\end{itemize}
\begin{lstlisting}
TeamThreadRange(TeamPolicy<...>::member_type, 
                begin, end)
\end{lstlisting}
\vspace{-0.4cm}
\begin{itemize}\itemsep1pt \parskip0pt \parsep0pt
\item for each calling team split the range \lstinline|[begin,end)| over threads in the team
\item must be nested in call with \lstinline|TeamPolicy|
\item team collective
\end{itemize}        
\begin{lstlisting}
ThreadVectorRange(TeamPolicy<...>::member_type, 
                begin, end)
\end{lstlisting}
\vspace{-0.4cm}
\begin{itemize}\itemsep1pt \parskip0pt \parsep0pt
\item for each calling thread split the range \lstinline|[begin,end)| over vector lanes
\item must be nested in call with \lstinline|TeamPolicy|
\end{itemize}        
\end{small}
      \end{minipage}
      };
      \node[fancytitle, right=10pt] at (ExecutionPolicies.north west) {Execution Policies};   
      \end{tikzpicture}
  }

  {
  \begin{tikzpicture}[remember picture,overlay,trim left]
     \node[HighLightBox,yshift=-24.75cm,xshift=0.265\paperwidth] at (current page.north west) (ExecutionSpace)
     {\begin{minipage}[h]{0.4\paperwidth}
     \begin{small}
\begin{lstlisting}
Serial
\end{lstlisting}
\vspace{-0.4cm}
\begin{itemize}\itemsep1pt \parskip0pt \parsep0pt
\item basic serial execution
\end{itemize}
\vspace{-0.4cm}
\begin{lstlisting}
OpenMP
\end{lstlisting}
\vspace{-0.4cm}
\begin{itemize}\itemsep1pt \parskip0pt \parsep0pt
\item using OpenMP 3.1 as backend
\end{itemize}
\vspace{-0.4cm}
\begin{lstlisting}
Threads
\end{lstlisting}
\vspace{-0.4cm}
\begin{itemize}\itemsep1pt \parskip0pt \parsep0pt
\item using Pthreads on Unix based systems
\end{itemize}
\vspace{-0.4cm}
\begin{lstlisting}
Cuda
\end{lstlisting}
\vspace{-0.4cm}
\begin{itemize}\itemsep1pt \parskip0pt \parsep0pt
\item using NVIDIAs Cuda 6.5 or higher as backend
\end{itemize}
\vspace{-0.4cm}
\end{small}
      \end{minipage}
      };
      \node[fancytitle, right=10pt] at (ExecutionSpace.north west) {Execution Spaces};   
   \end{tikzpicture}
  }
\end{minipage}
\end{document}
